\documentclass[10pt]{article}
\usepackage[a4paper,margin=0.6in]{geometry}
\usepackage{amsmath,amssymb,amsfonts,mathtools}
\allowdisplaybreaks
\setlength{\jot}{5pt}
\numberwithin{equation}{section}

\title{Simplification of the finite terms ($\Lambda\to0$, $\vee\to\infty$): detailed note}
\author{}
\date{}

\begin{document}
\maketitle
\tableofcontents

\section{Summary}

The input is the list of finite terms in \texttt{Limiting\_Cases.tex}: Group I (Rows II--V, with c.c.), Group II (Rows I, IV) and Group III (Rows I--III, V, VI, with c.c.). The results of this note are:

\begin{enumerate}
\item \textbf{Group I, Row II} (momentum space). The colour factor is antisymmetric under $x\leftrightarrow y$, which removes two of the three logarithmic structures of $R^{i}$. The remaining kernel has a single logarithm, Eq.~\eqref{eq:calR}, and is exactly the Fourier transform of the coordinate-space kernel $K^{ij}\log[(x-z)^{2}/(x-w)^{2}]$ that appears in the other rows, Eq.~\eqref{eq:RowIIcoord}.
\item \textbf{Group I, Row III.} The logarithms of $(x-w)^{2}-(x-z)^{2}$ cancel completely, together with all poles at $(x-w)^{2}=(x-z)^{2}$. The 22 terms reduce to three structures, Eq.~\eqref{eq:RowIIIfinal}. One constant $i\sigma\pi$ survives; with the prefactor $i$ it is a real contribution.
\item \textbf{The $K\log$ sector.} Six contributions (I.II, I.III, I.IV, II.I first part, III.I, III.II) have the same kernel. Their colour factors add up to commutators of $\rho$ only, Eq.~\eqref{eq:Omegatot}. For commuting $\rho$ the whole sector cancels.
\item \textbf{Group II, Row I, second part} ($2a$--$2f$). The six pieces are a perfect square of one two-gluon amplitude, Eq.~\eqref{eq:square}, with the closed form \eqref{eq:KIIclosed}.
\item \textbf{Group III, Rows III (first part), V, VI} add up to one term whose colour factor, for commuting $\rho$, is a product of three LO factors, Eq.~\eqref{eq:C3}.
\item \textbf{Group I Row V and Group III Row III (second parts)} are written with the same three building blocks $E,G,\phi$ as the square of item 4; in Row V the two pieces $N_{2a}$ and $N_{2b\,\rm rest}$ merge.
\item \textbf{Longitudinal integrals} that contain only a phase are done in closed form (Group I Row V first part, Group II Row IV, Group III).
\item \textbf{Reality.} Every group is real. The only complex kernel is the $i\sigma\pi$ of item 2.
\end{enumerate}
The final list of what is left is in Section~\ref{sec:final}; assumptions and open points are in Section~\ref{sec:open}. Every algebraic statement was checked numerically (Appendix~\ref{app:checks}).

\section{Conventions and building blocks}\label{sec:conv}

\subsection{General}
\begin{align}
n^{i}(r)\equiv\frac{r^{i}}{r^{2}},\qquad M\cdot N\equiv M^{km}N^{km},\qquad (M^{T})^{km}\equiv M^{mk}.
\end{align}
Transverse vectors are written $a,b,c,\dots$; they should not be confused with colour indices in $U^{ab}$. ``$+{\rm c.c.}$'' is the hermitian conjugate for operator-valued $\rho$. For Group II, $N_{1b}=(N_{1a})^{*}$, so $N_{1a}+N_{1b}$ is also of the form $X+{\rm c.c.}$

\subsection{Kernels of one pair of vectors}
For a pair $(a,b)$ with $c\equiv a-b$, $A\equiv a^{2}$, $B\equiv b^{2}$, $c^{2}=A+B-2a\cdot b$:
\begin{align}
K^{ij}(a,b)&\equiv\frac{1}{c^{2}}\left\{\frac{\delta^{ij}}{2}-\frac{a^{i}a^{j}}{A}-\frac{b^{i}b^{j}}{B}+\frac{2\,(a\cdot b)\,a^{i}b^{j}}{AB}\right\}, \label{eq:K}\\
S^{ij}(a,b)&\equiv\frac{1}{c^{2}}\left\{\delta^{ij}-\frac{a^{i}a^{j}}{A}-\frac{b^{i}b^{j}}{B}+\frac{(A+B)\,a^{i}b^{j}}{AB}\right\}=K^{ij}+\frac{\delta^{ij}}{2c^{2}}+\frac{a^{i}b^{j}}{AB}, \label{eq:S}\\
T^{ij}(a,b)&\equiv\frac{1}{c^{2}}\left\{\frac{(c\cdot a)}{A}\,a^{i}a^{j}+\frac{(c\cdot b)}{B}\,b^{i}b^{j}-\frac{(a\cdot b)\,(A-B)}{AB}\,a^{i}b^{j}\right\}, \label{eq:T}
\end{align}
and
\begin{align}
E^{km}(a,b)\equiv a^{k}\left[\frac{c^{m}}{c^{2}}+\frac{b^{m}}{2B}\right],\qquad
G^{km}(a,b)\equiv\left[\frac{c^{k}}{c^{2}}-\frac{a^{k}}{2A}\right]b^{m},\qquad
\phi(a,b)\equiv\frac{B-A}{2\,c^{2}}, \label{eq:blocks}
\end{align}
with traces $e\equiv E^{kk}=\frac{a\cdot c}{c^{2}}+\frac{a\cdot b}{2B}$ and $g\equiv G^{kk}=\frac{b\cdot c}{c^{2}}-\frac{a\cdot b}{2A}$. Identities used repeatedly:
\begin{align}
c\cdot a=A-a\cdot b,\qquad c\cdot b=a\cdot b-B,\qquad A-B=c\cdot(a+b). \label{eq:ids}
\end{align}
Unless stated otherwise the pair is
\begin{align}
a=x-w,\qquad b=x-z,\qquad c=z-w .
\end{align}

\subsection{Colour objects}
\begin{align}
\mathcal{D}_{0}^{a}&\equiv\big[U^{ab'}(x')-U^{ab'}(w')\big]\rho^{b'}(x'), &
\mathcal{D}^{c}(u,v)&\equiv\big[U^{ce}(u)-U^{ce}(v)\big]\rho^{e}(v),\notag\\
R^{ab}&\equiv f^{cbd}\,U^{ac}(w)\,\rho^{d}(x)-f^{acd}\,U^{cb}(z)\,U^{de}(x)\,\rho^{e}(x), &
V^{be}&\equiv U^{cb}(z)\,U^{ce}(y),\notag\\
\mathcal{S}^{a}&\equiv f^{hbc}\,U^{ah}(w)\{\rho^{b}(x),\rho^{c}(y)\}-f^{abc}\,U^{bd}(x)\,U^{ce}(y)\{\rho^{d}(x),\rho^{e}(y)\}. \label{eq:colour}
\end{align}
$R'^{ab}$ is $R^{ab}$ with $(w,x)\to(w',x')$ and the same $z$. $\mathcal{S}^{a}$ is the colour bracket of Group I Row II.

\subsection{Special functions and plus prescriptions}
\begin{align}
E_{1}(z)=\int_{1}^{\infty}\frac{dt}{t}\,e^{-zt},\qquad {\rm Ein}(z)=\int_{0}^{1}\frac{dt}{t}\,(1-e^{-zt})=\gamma_{E}+\log z+E_{1}(z).
\end{align}
For real $s>0$: $E_{1}(is)=-{\rm Ci}(s)+i\,[{\rm Si}(s)-\tfrac{\pi}{2}]$, ${\rm Ein}(is)=\gamma_{E}+\log s-{\rm Ci}(s)+i\,{\rm Si}(s)$. The plus prescriptions are read as
\begin{align}
\left[\frac{1}{\xi}\right]_{+}f\equiv\frac{f(\xi)-f(0)}{\xi}\ \ (0<\xi<1),\qquad
\left[\frac{1}{\eta(1+\eta)}\right]_{+}f\equiv\frac{f(\eta)-f(0)}{\eta(1+\eta)}\ \ (0<\eta<\infty). \label{eq:plus}
\end{align}

\section{Group I}

\subsection{Row II}\label{sec:RowII}

\paragraph{Step 1: rename.} The integration momentum of Row II is called $k$ in the file, the same symbol as the external momentum. I call it $q$, and I assume the integral over $x'$ is implied:
\begin{align}
\frac{d^{3}N_{\rm I.II}}{d^{3}k}=-\frac{1}{(2\pi)^{3}}\,\frac{g^{4}}{64\pi^{6}}\,\frac{1}{k^{+}}\int_{w,w',x'}e^{-ik\cdot(w'-w)}\int d^{2}p\,d^{2}q\int_{x,y}e^{iq\cdot(w-y)}\,e^{ip\cdot(y-x)}\;n^{i}(x'-w')\,\frac{R^{i}(q,p)}{p^{2}q^{2}}\;\mathcal{D}_{0}^{a}\,\mathcal{S}^{a}+{\rm c.c.}
\end{align}
\begin{align}
R^{i}(q,p)=-2q^{i}\log\frac{p^{2}}{q^{2}}+\frac{(2p^{2}-2\,q\cdot p+q^{2})\,q^{i}+q^{2}p^{i}}{(q-p)^{2}}\,\log\frac{q^{2}}{(q-p)^{2}} .
\end{align}

\paragraph{Step 2: the colour factor is antisymmetric under $x\leftrightarrow y$.} Exchanging $x\leftrightarrow y$ in $\mathcal{S}^{a}$ and relabelling $b\leftrightarrow c$ ($d\leftrightarrow e$) gives $f^{hcb}=-f^{hbc}$ while the anticommutators are symmetric, hence
\begin{align}
\mathcal{S}^{a}\big|_{x\leftrightarrow y}=-\,\mathcal{S}^{a}
\end{align}
exactly, also for non-commuting $\rho$. Nothing else in the integrand depends on $x,y$ except the phases, so only the part of the kernel that is antisymmetric under $x\leftrightarrow y$ contributes.

\paragraph{Step 3: action of $x\leftrightarrow y$ on the momenta.} With $l\equiv p-q$,
\begin{align}
e^{iq\cdot(w-y)}e^{ip\cdot(y-x)}\;\xrightarrow{x\leftrightarrow y}\;e^{iq\cdot(w-x)}e^{ip\cdot(x-y)}=e^{iq\cdot(w-y)}\,e^{i(q-p)\cdot(y-x)},
\end{align}
i.e.\ $p\to q-p$ at fixed $q$, which is $p^{2}\leftrightarrow l^{2}$. A term of the kernel that is symmetric under $p^{2}\leftrightarrow l^{2}$ at fixed $q$ drops out.

\paragraph{Step 4: split the kernel.} Using $2p^{2}-2q\cdot p+q^{2}=p^{2}+l^{2}$ and $p=q+l$,
\begin{align}
\frac{R^{i}}{p^{2}q^{2}}&=-\frac{2q^{i}}{p^{2}q^{2}}\log\frac{p^{2}}{q^{2}}+\left[\frac{q^{i}}{q^{2}l^{2}}+\frac{q^{i}}{p^{2}q^{2}}+\frac{p^{i}}{p^{2}l^{2}}\right]\log\frac{q^{2}}{l^{2}}
=\mathcal{R}^{i}(q,p)+2\,\mathcal{E}^{i}(q,p),
\end{align}
\begin{align}
\boxed{\;\mathcal{R}^{i}(q,p)\equiv\log\frac{q^{2}}{(p-q)^{2}}\left[\frac{p^{i}}{p^{2}\,(p-q)^{2}}-\frac{q^{i}}{q^{2}\,(p-q)^{2}}+\frac{q^{i}}{p^{2}q^{2}}\right]\;},\qquad
\mathcal{E}^{i}\equiv\frac{q^{i}}{q^{2}}\left[\frac{1}{l^{2}}\log\frac{q^{2}}{l^{2}}-\frac{1}{p^{2}}\log\frac{p^{2}}{q^{2}}\right]. \label{eq:calR}
\end{align}
$\mathcal{E}^{i}$ is symmetric under $p^{2}\leftrightarrow l^{2}$ and drops out by Step 3. It contains the whole term $-2q^{i}\log(p^{2}/q^{2})$ and the non-integrable $l\to0$ behaviour $\sim\log(l^{2})/l^{2}$ of the original kernel; $\mathcal{R}^{i}$ is only $\mathcal{O}(1/|l|)$ at $l\to0$ because $p^{i}/p^{2}-q^{i}/q^{2}=\mathcal{O}(l)$. So
\begin{align}
\frac{d^{3}N_{\rm I.II}}{d^{3}k}=-\frac{1}{(2\pi)^{3}}\,\frac{g^{4}}{64\pi^{6}}\,\frac{1}{k^{+}}\int_{w,w',x'}e^{-ik\cdot(w'-w)}\int d^{2}p\,d^{2}q\int_{x,y}e^{iq\cdot(w-y)}\,e^{ip\cdot(y-x)}\;n^{i}(x'-w')\,\mathcal{R}^{i}(q,p)\;\mathcal{D}_{0}^{a}\,\mathcal{S}^{a}+{\rm c.c.} \label{eq:RowIImom}
\end{align}

\paragraph{Step 5: coordinate space.} The Fourier transform of the kernel $K^{ij}\log(b^{2}/a^{2})$ is computed in Appendix~\ref{app:FT}:
\begin{align}
l^{j}\int d^{2}a\,d^{2}b\;e^{iq\cdot a+il\cdot b}\,K^{ij}(a,b)\,\log\frac{b^{2}}{a^{2}}=-2\pi^{2}\,l^{2}\,\mathcal{R}^{i}(q,q+l). \label{eq:FTresult}
\end{align}
Define $H^{i}\equiv\int_{z}n^{j}(y-z)\,K^{ij}(a,b)\log(B/A)$ with $a=x-w$, $b=x-z$, and $\rho\equiv y-x$ so that $y-z=\rho+b$. Its transform is
\begin{align}
\int d^{2}a\,d^{2}\rho\;e^{iq\cdot a-il\cdot\rho}\,H^{i}
=\underbrace{\int d^{2}r\,e^{-il\cdot r}\,n^{j}(r)}_{-2\pi i\,l^{j}/l^{2}}\;\int d^{2}a\,d^{2}b\;e^{iq\cdot a+il\cdot b}K^{ij}\log\frac{b^{2}}{a^{2}}=4\pi^{3}i\;\mathcal{R}^{i}(q,q+l).
\end{align}
Since $e^{iq\cdot(w-y)}e^{ip\cdot(y-x)}=e^{-iq\cdot a}e^{il\cdot\rho}$, inverting gives
\begin{align}
\int d^{2}p\,d^{2}q\;e^{iq\cdot(w-y)}e^{ip\cdot(y-x)}\,\mathcal{R}^{i}(q,p)=\frac{(2\pi)^{4}}{4\pi^{3}i}\,H^{i}=-4\pi i\,H^{i},
\end{align}
and \eqref{eq:RowIImom} becomes
\begin{align}
\boxed{\;\frac{d^{3}N_{\rm I.II}}{d^{3}k}=\frac{1}{(2\pi)^{3}}\,\frac{ig^{4}}{8\pi^{5}}\,\frac{1}{k^{+}}\int_{w,w',x,y,z,x'}e^{-ik\cdot(w'-w)}\;n^{i}(x'-w')\,n^{j}(y-z)\,K^{ij}\,\log\frac{(x-z)^{2}}{(x-w)^{2}}\;\;\tfrac12\,\mathcal{D}_{0}^{a}\,\mathcal{S}^{a}+{\rm c.c.}\;} \label{eq:RowIIcoord}
\end{align}
Row II is therefore one more term of the $K\log$ type (Section~\ref{sec:Ksector}).

\subsection{Row III}\label{sec:RowIII}

All terms of $N_{1},N_{2},N_{3a},N_{3b}$ carry the prefactor $ig^{4}/8\pi^{5}k^{+}$ (some with an extra $\tfrac12$) and the same colour factor
\begin{align}
\mathcal{W}_{\rm III}\equiv\mathcal{D}_{0}^{a}\,f^{abd}\,U^{de}(x)\big[U^{bc}(z)-U^{bc}(y)\big]\rho^{e}(x)\,\rho^{c}(y),
\end{align}
which does not depend on $w$. With $u^{i}\equiv n^{i}(x'-w')$, $v^{j}\equiv n^{j}(y-z)$ write
\begin{align}
\frac{d^{3}N_{\rm I.III}}{d^{3}k}=\frac{1}{(2\pi)^{3}}\,\frac{i g^{4}}{8\pi^{5}}\,\frac{1}{k^{+}}\int_{w,w',x,y,z,x'}e^{-ik\cdot(w'-w)}\;u^{i}\,v^{j}\;\mathcal{R}_{\rm III}^{ij}\;\mathcal{W}_{\rm III}\;+\;{\rm c.c.} \label{eq:defR}
\end{align}

\paragraph{Step 1: the four pieces.} Reading off the file (the overall sign of $N_{2}$ is $-$):
\begin{align}
\mathcal{R}^{ij}_{1}&=\frac{\delta^{ij}}{2}\,\frac{A+B}{(B-A)^{2}}\left\{\log\frac{B}{B-A}+\log\frac{A-B}{A}\right\}-\frac{a^{i}\,b^{j}}{A\,B}\,\log\frac{-B}{A},\notag\\
\mathcal{R}^{ij}_{2}&=-\,M^{ij}\,\log\frac{-B}{A-B},\notag\\
\mathcal{R}^{ij}_{3a}&=-\,M^{ij}\,\log\frac{A-B}{A}
+\mathcal{A}_{0}\left[\frac{c^{i}a^{j}}{c^{2}}+\frac{b^{i}c^{j}}{c^{2}}-\frac{(a\cdot b)\,c^{i}b^{j}}{c^{2}\,B}\right],\notag\\
\mathcal{R}^{ij}_{3b}&=-\,\mathcal{A}_{0}\,\frac{(a\cdot b)\,a^{i}c^{j}}{c^{2}\,A}
-\frac12\,\log\frac{c^{2}}{A}\left[\frac{c^{i}b^{j}}{c^{2}B}-\frac{a^{i}c^{j}}{c^{2}A}+\frac{\delta^{ij}}{c^{2}}\right],
\end{align}
where the same tensor appears in $N_{2}$ (six terms) and in the six logarithmic terms of $N_{3a}$ (with opposite sign):
\begin{align}
M^{ij}\equiv\frac{1}{c^{2}}\Bigg\{\frac{c^{i}a^{j}+b^{i}c^{j}}{A-B}-\frac{1}{A-B}\left[\frac{B}{A}\,a^{i}c^{j}+\frac{A}{B}\,c^{i}b^{j}\right]
+\delta^{ij}\,\frac{B\,(c\cdot a)-A\,(c\cdot b)}{(A-B)^{2}}\Bigg\}.
\end{align}

\paragraph{Step 2: the logarithms of $(A-B)$ cancel.} With
$L_{a}\equiv\log\frac{B}{B-A}=\log\frac{-B}{A-B}$, $L_{b}\equiv\log\frac{A-B}{A}$, $L_{c}\equiv\log\frac{-B}{A}$:
for $A>B$ the argument of $L_{a}$ is negative and that of $L_{b}$ positive, for $A<B$ the reverse, and in both regions
\begin{align}
L_{a}+L_{b}=\log\left[\frac{B}{B-A}\cdot\frac{A-B}{A}\right]=L_{c}=\log\frac{B}{A}+i\,\sigma\,\pi ,\qquad\sigma=\pm1. \label{eq:logsum}
\end{align}
$\log|A-B|$ drops out. The remnant of the negative arguments is the constant $i\sigma\pi$ with the same sign in both regions ($\sigma=+1$ for the principal branch, $\sigma=-1$ for the opposite $i\epsilon$). Then $\mathcal{R}_{2}+\mathcal{R}_{3a}|_{\log}=-M^{ij}(L_{a}+L_{b})=-M^{ij}L_{c}$ and
\begin{align}
\mathcal{R}_{\rm III}^{ij}\Big|_{L_{c}}=\log\frac{-B}{A}\left\{\frac{\delta^{ij}}{2}\,\frac{A+B}{(A-B)^{2}}-\frac{a^{i}b^{j}}{AB}-M^{ij}\right\}.
\end{align}

\paragraph{Step 3: the poles at $A=B$ cancel.} For the $\delta^{ij}$ part, \eqref{eq:ids} gives $B\,(c\cdot a)-A\,(c\cdot b)=2AB-(a\cdot b)(A+B)$ and
\begin{align}
\frac{A+B}{2(A-B)^{2}}-\frac{2AB-(a\cdot b)(A+B)}{c^{2}(A-B)^{2}}
=\frac{\tfrac12(A+B)(A+B-2a\cdot b)-2AB+(a\cdot b)(A+B)}{c^{2}(A-B)^{2}}
=\frac{1}{2c^{2}} .
\end{align}
For the rest, with $c=a-b$: $c^{i}a^{j}+b^{i}c^{j}=a^{i}a^{j}-b^{i}b^{j}$ and
$\frac{B}{A}a^{i}c^{j}+\frac{A}{B}c^{i}b^{j}=\frac{B}{A}a^{i}a^{j}-\frac{A}{B}b^{i}b^{j}+a^{i}b^{j}\frac{A^{2}-B^{2}}{AB}$, so
\begin{align}
M^{ij}\Big|_{{\rm non-}\delta}
=\frac{1}{c^{2}}\left\{\frac{a^{i}a^{j}}{A}+\frac{b^{i}b^{j}}{B}-\frac{(A+B)\,a^{i}b^{j}}{AB}\right\}.
\end{align}
Adding $-a^{i}b^{j}/(AB)=-c^{2}a^{i}b^{j}/(c^{2}AB)$ and using $A+B-c^{2}=2a\cdot b$,
\begin{align}
\mathcal{R}_{\rm III}^{ij}\Big|_{L_{c}}=K^{ij}\,\log\frac{-B}{A},
\end{align}
with $K^{ij}$ of \eqref{eq:K}. Writing $c=a-b$ in the bracket of Group I Row IV (the same bracket appears in Group II Row I ($1a$) and Group III Rows I, II) shows it is the same tensor:
\begin{align}
\frac{(x-z)^{j}(z-w)^{i}}{(z-w)^{2}(x-z)^{2}}-\frac{(x-w)^{i}(z-w)^{j}}{(z-w)^{2}(x-w)^{2}}-\frac{(x-w)^{i}(x-z)^{j}}{(x-w)^{2}(x-z)^{2}}+\frac{\delta^{ij}}{2(z-w)^{2}}=K^{ij}.
\end{align}

\paragraph{Step 4: the $\mathcal{A}_{0}$ terms.}
\begin{align}
\mathcal{R}_{\rm III}^{ij}\Big|_{\mathcal{A}_{0}}=\frac{\mathcal{A}_{0}}{c^{2}}\left\{c^{i}a^{j}+b^{i}c^{j}-(a\cdot b)\left[\frac{c^{i}b^{j}}{B}+\frac{a^{i}c^{j}}{A}\right]\right\}=\mathcal{A}_{0}\,T^{ij},
\end{align}
with $T^{ij}$ of \eqref{eq:T}. With $\Delta\equiv\sqrt{AB-(a\cdot b)^{2}}$ (twice the area of the triangle $x,z,w$), $\arctan[(a\cdot b)/\Delta]=\tfrac\pi2-\theta_{x}$ and $\arctan[(B-a\cdot b)/\Delta]=\tfrac\pi2-\theta_{z}$ where $\theta_{x},\theta_{z}$ are the interior angles at $x$ and $z$. Hence
\begin{align}
\mathcal{A}_{0}=\frac{\pi-\theta_{x}-\theta_{z}}{\Delta}=\frac{\theta_{w}}{\sqrt{(x-w)^{2}(x-z)^{2}-[(x-w)\cdot(x-z)]^{2}}},\qquad
\cos\theta_{w}=\frac{(x-w)\cdot(z-w)}{|x-w|\,|z-w|},
\end{align}
the interior angle at $w$ divided by twice the area: real, no logarithm.

\paragraph{Step 5: the $\log[(z-w)^{2}/(x-w)^{2}]$ terms} have a positive argument and give $-\tfrac12\log\frac{c^{2}}{A}\,S^{ij}$ with $S^{ij}$ of \eqref{eq:S}.

\paragraph{Step 6: result.}
\begin{align}
\boxed{\;\mathcal{R}_{\rm III}^{ij}=K^{ij}\left\{\log\frac{(x-z)^{2}}{(x-w)^{2}}+i\,\sigma\,\pi\right\}
-\frac12\,S^{ij}\,\log\frac{(z-w)^{2}}{(x-w)^{2}}+\mathcal{A}_{0}\,T^{ij}\;} \label{eq:RowIIIfinal}
\end{align}

\subsection{Row IV}
Using $U^{bd}(z)U^{be}(y)=V^{de}$ the colour bracket of Row IV is $V^{be}\rho^{e}(y)\,R^{ab}$ (first term: $V^{de}\rho^{e}(y)f^{cdc'}U^{ac}(w)\rho^{c'}(x)$; second term: $U^{cb}(z)U^{gb}(z)=\delta^{cg}$). So
\begin{align}
\frac{d^{3}N_{\rm I.IV}}{d^{3}k}=\frac{1}{(2\pi)^{3}}\,\frac{ig^{4}}{8\pi^{5}}\,\frac{1}{k^{+}}\int e^{-ik\cdot(w'-w)}\,u^{i}v^{j}K^{ij}\log\frac{(x-z)^{2}}{(x-w)^{2}}\;\mathcal{D}_{0}^{a}\,V^{be}\rho^{e}(y)\,R^{ab}+{\rm c.c.}
\end{align}

\subsection{Row V, first part}
With $\sigma\equiv k\cdot(z-x)$ and \eqref{eq:plus}:
\begin{align}
\int_{0}^{1}d\xi\,e^{-i\xi\sigma}=\frac{1-e^{-i\sigma}}{i\sigma}\equiv h_{0}(\sigma),\qquad
\int_{0}^{1}d\xi\left[\frac1\xi\right]_{+}e^{\mp i\xi\sigma}=-{\rm Ein}(\pm i\sigma).
\end{align}
The prefactors are $\frac{1}{k^{+2}}\int_{0}^{k^{+}}\frac{dp^{+}}{2\pi}=\frac{1}{2\pi k^{+}}\int_{0}^{1}d\xi$ and $\frac{1}{k^{+}}\int_{0}^{k^{+}}\frac{dq^{+}}{2\pi}[\frac{1}{q^{+}}]_{+}=\frac{1}{2\pi k^{+}}\int_{0}^{1}d\zeta\,[\frac1\zeta]_{+}$. Then
\begin{align}
\frac{d^{3}N_{\rm I.V(1)}}{d^{3}k}&=\frac{1}{(2\pi)^{3}}\,\frac{ig^{4}}{32\pi^{5}}\,\frac{1}{k^{+}}\int_{x,z,y,x',y',w'}\Big\{e^{-ik\cdot(w'-z)}\Big[h_{0}(\sigma)\;n(x'-w')\!\cdot\!n(z-x)\;\,n(y-z)\!\cdot\!n(y'-x)\notag\\
&\hspace{2cm}+{\rm Ein}(i\sigma)\;n(y'-x)\!\cdot\!n(z-x)\;\,n(x'-w')\!\cdot\!n(y-z)\Big]\notag\\
&\quad+e^{-ik\cdot(w'-x)}\,{\rm Ein}(-i\sigma)\;n(x'-w')\!\cdot\!n(y'-x)\;\,n(z-x)\!\cdot\!n(y-z)\Big\}\;\mathcal{D}_{0}^{a}\notag\\
&\quad\times\Big[f^{adc}\,U^{db}(x)U^{ce}(z)\{\rho^{e}(y),\rho^{b}(y')\}-f^{abc}\,U^{cd}(y)U^{be}(y')\{\rho^{d}(y),\rho^{e}(y')\}\Big]+{\rm c.c.} \label{eq:RowV1}
\end{align}

\subsection{Row V, second part: $N_{2a}+N_{2b\,\rm rest}$}
Use the blocks \eqref{eq:blocks} at $\bar a\equiv y-x$, $\bar b\equiv y-z$, $\bar c=\bar a-\bar b=z-x$ (bars on $E,G,\phi,e,g$), and $\xi=p^{+}/k^{+}$.

\paragraph{Step 1: the bracket of $N_{2b\,\rm rest}$.} Term by term ($x-z=-\bar c$):
\begin{align}
\frac{\delta^{im}d_{\perp}}{2(x-z)^{2}}\big[\bar a^{2}-\bar b^{2}\big]\xi(1-\xi)&=-d_{\perp}\,\xi(1-\xi)\,\bar\phi\,\delta^{im},\notag\\
\frac{(y-x)^{i}(x-z)^{m}}{(x-z)^{2}}-\frac{(y-x)^{i}(y-z)^{m}}{2(y-z)^{2}}&=-\bar E^{im},\notag\\
\frac{(y-z)\cdot(x-z)}{(x-z)^{2}}+\frac{(y-x)\cdot(y-z)}{2(y-x)^{2}}-\frac{\bar a^{2}-\bar b^{2}}{2(x-z)^{2}}&=-\bar g+\bar\phi,\notag\\
\frac{(y-z)^{i}(x-z)^{m}}{(x-z)^{2}}+\frac{(y-x)^{m}(y-z)^{i}}{2(y-x)^{2}}&=-\bar G^{mi},\notag\\
\frac{(y-x)\cdot(x-z)}{(x-z)^{2}}-\frac{(y-x)\cdot(y-z)}{2(y-z)^{2}}-\frac{\bar a^{2}-\bar b^{2}}{2(x-z)^{2}}&=-\bar e+\bar\phi,\notag\\
\frac{(y-x)^{m}(x-z)^{i}}{(x-z)^{2}}-\frac{(y-x)^{m}(y-z)^{i}}{2(y-z)^{2}}+\frac{(y-z)^{m}(x-z)^{i}}{(x-z)^{2}}+\frac{(y-x)^{i}(y-z)^{m}}{2(y-x)^{2}}&=-\bar E^{mi}-\bar G^{im}.
\end{align}
The bracket is therefore
\begin{align}
\mathfrak{B}_{V}^{im}(\xi)&=\delta^{im}\Big[\bar\phi\,\big(1-d_{\perp}\,\xi(1-\xi)\big)-\xi\,\bar e-(1-\xi)\,\bar g\Big]\notag\\
&\quad+\left(\left[\frac{1}{1-\xi}\right]_{+}-1\right)\bar E^{im}+\left(\left[\frac{1}{\xi}\right]_{+}-1\right)\bar G^{mi}+\bar E^{mi}+\bar G^{im}. \label{eq:BV}
\end{align}

\paragraph{Step 2: $N_{2a}$ has the same bracket.} Its first integral is \eqref{eq:BV} without the $\bar E^{im}$ term. In its second integral put $q^{+}=k^{+}-p^{+}$: the phase $e^{-ik\cdot(w'-x)}e^{-ik\cdot(x-z)q^{+}/k^{+}}$ becomes $e^{-ik\cdot(w'-z)}e^{-ik\cdot(z-x)p^{+}/k^{+}}$, the denominator becomes $p^{+}\bar a^{2}+(k^{+}-p^{+})\bar b^{2}$, and $\int^{0}_{k^{+}}dq^{+}=-\int_{0}^{k^{+}}dq^{+}$ (limits as written in the file) gives $+([\tfrac{1}{1-\xi}]_{+}-1)\bar E^{im}$. So $N_{2a}$ and $N_{2b\,\rm rest}$ share the kernel $\mathfrak{B}_{V}^{im}/[\xi\bar a^{2}+(1-\xi)\bar b^{2}]$.

\paragraph{Step 3: add the colour factors.}
$\big[-f^{adc}f^{fbe}U^{db}(x)U^{ce}(z)+N_{c}U^{af}(x)\big]+N_{c}\big[U^{af}(y)-U^{af}(x)\big]$: the $U(x)$ terms cancel.
\begin{align}
\frac{d^{3}N_{\rm I.V(2)}}{d^{3}k}&=\frac{1}{(2\pi)^{3}}\,\frac{g^{4}}{16\pi^{5}}\,\frac{1}{k^{+}}\int_{0}^{1}d\xi\int_{x,y,z,x',w'}\Bigg\{
e^{-ik\cdot(w'-z)}\,e^{-i\xi\,k\cdot(z-x)}\;\frac{n^{i}(x'-w')\,\mathfrak{B}_{V}^{im}(\xi)\,n^{m}(z-x)}{\xi\,(y-x)^{2}+(1-\xi)\,(y-z)^{2}}\notag\\
&\hspace{1.5cm}\times\big[U^{ab'}(x')-U^{ab'}(w')\big]\Big[N_{c}\,U^{af}(y)-f^{adc}f^{fbe}\,U^{db}(x)\,U^{ce}(z)\Big]\notag\\
&-e^{-ik\cdot(w'-x)}\,\big[\xi(1-\xi)-2\big]\;\frac{n(x'-w')\!\cdot\!(y-x)\;\;n(y-z)\!\cdot\!n(y-x)}{(x-z)^{2}}\notag\\
&\hspace{1.5cm}\times N_{c}\big[U^{ab'}(x')-U^{ab'}(w')\big]\big[U^{af}(y)-U^{af}(x)\big]\Bigg\}\rho^{b'}(x')\rho^{f}(y)+{\rm c.c.} \label{eq:RowV2}
\end{align}
The colour factor of the first term is that of a gluon at $y$ splitting into gluons at $x$ and $z$. At $z\to x$ it becomes $N_{c}[U(y)-U(x)]$, because $f^{adc}f^{fbe}U^{db}(x)U^{ce}(x)=N_{c}U^{af}(x)$, which is the colour factor of the subtraction, and $\xi(1-\xi)-2$ is the plus-prescription form of $\tfrac{\xi}{1-\xi}+\tfrac{1-\xi}{\xi}+\xi(1-\xi)$.

\section{Group II}

\subsection{Row I, first part}
With $u^{i}=n^{i}(x'-w')$, $v^{j}=n^{j}(y-z)$,
\begin{align}
\frac{d^{3}N_{1a}}{d^{3}k}&=-\frac{1}{(2\pi)^{3}}\,\frac{ig^{4}}{16\pi^{5}}\,\frac{1}{k^{+}}\int e^{-ik\cdot(w'-w)}\,u^{i}v^{j}K^{ij}\,\log\frac{(x-z)^{2}}{(x-w)^{2}}\notag\\
&\quad\times\Big[U^{ac'}(w')\{\rho^{c'}(x'),\rho^{b}(y)\}-U^{ad'}(x')\,V^{be'}\{\rho^{d'}(x'),\rho^{e'}(y)\}\Big]R^{ab},\qquad N_{1b}=(N_{1a})^{*}. \label{eq:GII1}
\end{align}

\subsection{Row I, second part ($2a$--$2f$)}\label{sec:square}

\paragraph{Step 1: notation and colour.}
\begin{align}
&a=x-w,\quad b=x-z,\quad c=z-w,\qquad a'=x'-w',\quad b'=x'-z,\quad c'=z-w'=a'-b',\notag\\
&\ell\equiv\log\frac{B}{A},\qquad \ell'\equiv\log\frac{B'}{A'},\qquad D\equiv A'B-AB' .
\end{align}
A prime on $E,G,\phi,e,g$ means evaluation at $(a',b')$. Using $U^{c'b}(z)U^{cb}(z)=\delta^{c'c}$ and $f^{acd'}f^{acd}=N_{c}\delta^{d'd}$, the four-term colour bracket common to all six pieces is
\begin{align}
f^{c'bd'}f^{cbd}\,U^{ac'}(w')U^{ac}(w)\rho^{d'}(x')\rho^{d}(x)-\dots+N_{c}\,U^{de'}(x')U^{de}(x)\rho^{e'}(x')\rho^{e}(x)=R'^{ab}R^{ab}.
\end{align}
So
\begin{align}
\frac{d^{3}N_{2a+\dots+2f}}{d^{3}k}=\frac{1}{(2\pi)^{3}}\,\frac{g^{4}}{8\pi^{5}}\,\frac{1}{k^{+}}\int_{z,w,w',x,x'}e^{-ik\cdot(w'-w)}\;\mathcal{K}_{\rm II}\;R'^{ab}R^{ab}. \label{eq:GIIdef}
\end{align}

\paragraph{Step 2: the brackets are products of the blocks.}
\begin{align}
P&\equiv[\dots]_{2a}=E\cdot G'+G\cdot E', & X&\equiv[\dots]_{2b}=E\cdot E', & Y&\equiv[\dots]_{2c}=G\cdot G',\notag\\
Z&\equiv[\dots]_{2d}=\phi'\,e+\phi\,e', & V_{e}&\equiv[\dots]_{2e}=\phi'\,g+\phi\,g', & &\frac{d_{\perp}}{4}\frac{(B-A)(B'-A')}{c^{2}c'^{2}}=d_{\perp}\,\phi\,\phi' .
\end{align}
For example $E\cdot E'=(a\cdot a')\big[\tfrac{c}{c^{2}}+\tfrac{b}{2B}\big]\!\cdot\!\big[\tfrac{c'}{c'^{2}}+\tfrac{b'}{2B'}\big]$ is the bracket of $2b$; $E\cdot G'$ gives terms 1--4 of $2a$ and $G\cdot E'$ terms 5--8.

\paragraph{Step 3: coefficients of the logarithms} (in units of the prefactor of \eqref{eq:GIIdef}):
\begin{align}
2a&:\; \frac{P}{D}\,(\ell-\ell'), &
2d&:\; \frac{Z}{D}\left[\frac{B}{B-A}\,\ell-\frac{B'}{B'-A'}\,\ell'\right],\notag\\
2b&:\; \frac{X}{D}\left[-\frac{B}{A}\,\ell+\frac{B'}{A'}\,\ell'\right], &
2e&:\; \frac{V_{e}}{D}\left[-\frac{A}{B-A}\,\ell+\frac{A'}{B'-A'}\,\ell'\right],\notag\\
2c&:\; \frac{Y}{D}\left[-\frac{A}{B}\,\ell+\frac{A'}{B'}\,\ell'\right], &
2f&:\; -d_{\perp}\phi\phi'\left\{\frac{1}{D}\left[\frac{AB\,\ell}{(B-A)^{2}}-\frac{A'B'\,\ell'}{(B'-A')^{2}}\right]+\frac{1}{(B-A)(B'-A')}\right\}.
\end{align}

\paragraph{Step 4: each combination is one integral.} With $\Delta(\eta)\equiv(\eta B+A)(\eta B'+A')$, by partial fractions:
\begin{align}
2a:\quad\int_{0}^{\infty}\frac{d\eta}{\Delta}&=\frac{\ell-\ell'}{D},\notag\\
2b:\quad\int_{0}^{\infty}\frac{d\eta}{\eta}\left[\frac{1}{\Delta}-\frac{\theta(1-\eta)}{AA'}\right]&=\frac{1}{D}\left[-\frac{B}{A}\ell+\frac{B'}{A'}\ell'\right],\notag\\
2c:\quad\int_{0}^{\infty}d\eta\left[\frac{\eta}{\Delta}-\frac{\theta(\eta-1)}{BB'\,\eta}\right]&=\frac{1}{D}\left[-\frac{A}{B}\ell+\frac{A'}{B'}\ell'\right],\notag\\
2d:\quad\int_{0}^{\infty}\frac{d\eta}{(1+\eta)\,\Delta}&=\frac{1}{D}\left[\frac{B\,\ell}{B-A}-\frac{B'\,\ell'}{B'-A'}\right],\notag\\
2e:\quad\int_{0}^{\infty}\frac{\eta\,d\eta}{(1+\eta)\,\Delta}&=\frac{1}{D}\left[-\frac{A\,\ell}{B-A}+\frac{A'\,\ell'}{B'-A'}\right],\notag\\
2f:\quad\int_{0}^{\infty}\frac{\eta\,d\eta}{(1+\eta)^{2}\,\Delta}&=-\frac{1}{D}\left[\frac{AB\,\ell}{(B-A)^{2}}-\frac{A'B'\,\ell'}{(B'-A')^{2}}\right]-\frac{1}{(B-A)(B'-A')} .
\end{align}
The two subtractions are the endpoint subtractions at $\eta\to0$ (the $\Lambda\to0$ side) and $\eta\to\infty$ (the $\vee\to\infty$ side).

\paragraph{Step 5: perfect square.} Adding the six pieces,
\begin{align}
\mathcal{K}_{\rm II}=\int_{0}^{\infty}d\eta\;\frac{\mathcal{N}(\eta)}{\Delta(\eta)}\Bigg|_{\rm subtr.},\qquad
\mathcal{N}(\eta)=\frac{X}{\eta}+P+\eta\,Y+\frac{Z+\eta V_{e}}{1+\eta}+\frac{d_{\perp}\,\phi\phi'\,\eta}{(1+\eta)^{2}},
\end{align}
and by Step 2 the numerator factorises ($\delta^{kk}=d_{\perp}$):
\begin{align}
\mathcal{N}(\eta)=\Phi^{km}(a,b;\eta)\,\Phi^{km}(a',b';\eta),\qquad
\Phi^{km}(a,b;\eta)\equiv\frac{1}{\sqrt{\eta}}\,E^{km}+\sqrt{\eta}\,G^{km}+\frac{\sqrt{\eta}}{1+\eta}\,\phi\,\delta^{km}. \label{eq:Phi}
\end{align}
In the original variables
\begin{align}
\Phi^{km}=\frac{(x-w)^{k}}{\sqrt{\eta}}\left[\frac{(z-w)^{m}}{(z-w)^{2}}+\frac{(x-z)^{m}}{2(x-z)^{2}}\right]
+\sqrt{\eta}\left[\frac{(z-w)^{k}}{(z-w)^{2}}-\frac{(x-w)^{k}}{2(x-w)^{2}}\right](x-z)^{m}
+\frac{\sqrt{\eta}}{1+\eta}\,\frac{(x-z)^{2}-(x-w)^{2}}{2\,(z-w)^{2}}\,\delta^{km}.
\end{align}
Hence
\begin{align}
\boxed{\;\frac{d^{3}N_{2a+\dots+2f}}{d^{3}k}=\frac{1}{(2\pi)^{3}}\,\frac{g^{4}}{8\pi^{5}}\,\frac{1}{k^{+}}\int_{z,w,w',x,x'}e^{-ik\cdot(w'-w)}\int_{0}^{\infty}d\eta\;
\left[\frac{\Phi'^{km}\,R'^{ab}}{\eta B'+A'}\right]\left[\frac{\Phi^{km}\,R^{ab}}{\eta B+A}\right]\Bigg|_{\rm subtr.}\;} \label{eq:square}
\end{align}
``subtr.'': in the $E\cdot E'/\eta$ term $1/\Delta\to1/\Delta-\theta(1-\eta)/(AA')$; in the $\eta\,G\cdot G'$ term $\eta/\Delta\to\eta/\Delta-\theta(\eta-1)/(BB'\eta)$.

\paragraph{Closed form.} The coefficient of each logarithm is the residue of the integrand:
\begin{align}
\boxed{\;\mathcal{K}_{\rm II}=\frac{1}{D}\left\{\mathcal{N}\!\left(-\tfrac{A}{B}\right)\log\frac{B}{A}-\mathcal{N}\!\left(-\tfrac{A'}{B'}\right)\log\frac{B'}{A'}\right\}-\frac{d_{\perp}}{4\,c^{2}c'^{2}}\;} \label{eq:KIIclosed}
\end{align}
\begin{align}
\mathcal{N}\!\left(-\tfrac{A}{B}\right)=P-\frac{B}{A}X-\frac{A}{B}Y+\frac{BZ-AV_{e}}{B-A}-\frac{d_{\perp}\,\phi\phi'\,AB}{(B-A)^{2}},
\end{align}
and $\mathcal{N}(-A'/B')$ is the same with $A\to A'$, $B\to B'$ in the explicit coefficients.

\paragraph{Consequences.}
\begin{itemize}
\item At $D=0$ one has $A/B=A'/B'$: residues and logarithms coincide and the $1/D$ cancels. At $B=A$ the apparent poles are compensated by $\ell\simeq(B-A)/A$ and $\phi\propto(B-A)$.
\item $\mathcal{K}_{\rm II}$ is real and symmetric under (primed $\leftrightarrow$ unprimed). For commuting $\rho$ the unsubtracted integrand is $|\mathcal{M}|^{2}$ with $\mathcal{M}^{km,ab}(z,\eta)=\int_{w,x}e^{ik\cdot w}\Phi^{km}R^{ab}/(\eta B+A)$.
\item All logarithms have positive arguments.
\end{itemize}

\subsection{Row IV}
With $r\equiv y'-y$, $s\equiv k\cdot r$, $t=1/\xi$:
\begin{align}
\mathcal{H}(s)\equiv\int_{0}^{1}d\xi\,\frac{1-\xi}{\xi}\,e^{-is/\xi}=\int_{1}^{\infty}dt\left(\frac1t-\frac{1}{t^{2}}\right)e^{-ist}=(1+is)\,E_{1}(is)-e^{-is}.
\end{align}
For $s>0$: ${\rm Re}\,\mathcal{H}=-{\rm Ci}(s)-s[{\rm Si}(s)-\tfrac{\pi}{2}]-\cos s$, ${\rm Im}\,\mathcal{H}={\rm Si}(s)-\tfrac{\pi}{2}-s\,{\rm Ci}(s)+\sin s$, and $\mathcal{H}(-s)=\mathcal{H}(s)^{*}$. The bracket is the reflection tensor $\Pi^{ij}(r)\equiv\delta^{ij}-2r^{i}r^{j}/r^{2}$:
\begin{align}
\frac{d^{3}N_{\rm II.IV}}{d^{3}k}=\frac{1}{(2\pi)^{3}}\,\frac{g^{4}N_{c}}{8\pi^{4}}\,\frac{1}{k^{+}}\int_{x,y,x',y'}\mathcal{H}\big(k\cdot(y'-y)\big)\;n^{i}(x-y)\,\Pi^{ij}(y'-y)\,n^{j}(x'-y')\;
\mathcal{D}^{c}(y',x')\,\mathcal{D}^{c}(y,x). \label{eq:GIIIV}
\end{align}
The kernel is symmetric under $(x,y)\leftrightarrow(x',y')$ while $s\to-s$; for commuting $\rho$ only ${\rm Re}\,\mathcal{H}$ survives. $\mathcal{H}\simeq-\log|s|$ at $s\to0$ (integrable).

\section{Group III}

\subsection{Rows I and II}
Both have the kernel $u^{i}v^{j}K^{ij}\log(B/A)$, the prefactor $-ig^{4}/8\pi^{5}k^{+}$ and the right factor $R^{ab}$. With $U^{c'd'}(y)U^{c'b}(z)=V^{bd'}$ the left factors add to
\begin{align}
\frac{d^{3}N_{\rm III.I+III.II}}{d^{3}k}&=-\frac{1}{(2\pi)^{3}}\,\frac{ig^{4}}{8\pi^{5}}\,\frac{1}{k^{+}}\int e^{-ik\cdot(w'-w)}\,u^{i}v^{j}K^{ij}\log\frac{(x-z)^{2}}{(x-w)^{2}}\notag\\
&\quad\times\Big\{U^{ab'}(x')V^{bd'}\{\rho^{b'}(x'),\rho^{d'}(y)\}-U^{ac'}(w')V^{be'}\rho^{c'}(x')\rho^{e'}(y)-U^{ab'}(x')\rho^{b}(y)\rho^{b'}(x')\Big\}R^{ab}+{\rm c.c.}
\end{align}

\subsection{Row III (first part) $+$ Row V $+$ Row VI}
These have identical kernels, measures and phases; prefactors $ig^{4}/8\pi^{4}$, $ig^{4}/4\pi^{4}$, $ig^{4}/4\pi^{4}$. Let $\mathcal{D}^{c'}_{q}\equiv[U^{c'e'}(q)-U^{c'e'}(x')]\rho^{e'}(x')$ with $q=y'$ or $z$.

\paragraph{Step 1: colour.} In units of $ig^{4}/8\pi^{4}$, Rows III(1) and VI have the left factor $f^{c'd'a}U^{d'b}(z)\mathcal{D}^{c'}_{q}$ and right factors that add to
\begin{align}
&U^{ac}(w)\{\rho^{c}(y),\rho^{b}(x)\}-U^{cb}(z)U^{ad}(y)U^{ce}(x)\{\rho^{d}(y),\rho^{e}(x)\}+2\,U^{cb}(z)U^{ad}(y)U^{ce}(x)\rho^{d}(y)\rho^{e}(x)-2\,U^{ad}(y)\rho^{d}(y)\rho^{b}(x)\notag\\
&=2\big[U^{ac}(w)-U^{ac}(y)\big]\rho^{c}(y)\rho^{b}(x)-U^{ac}(w)\big[\rho^{c}(y),\rho^{b}(x)\big]+U^{cb}(z)U^{ad}(y)U^{ce}(x)\big[\rho^{d}(y),\rho^{e}(x)\big].
\end{align}
Row V is $2f^{c'ca}\mathcal{D}^{c'}_{q}U^{ce}(x)[U^{ad}(y)-U^{ad}(w)]\rho^{e}(x)\rho^{d}(y)=-2f^{c'd'a}\mathcal{D}^{c'}_{q}U^{d'b}(x)[U^{ac}(w)-U^{ac}(y)]\rho^{b}(x)\rho^{c}(y)$. Adding, and using $U^{d'b}(z)U^{gb}(z)=\delta^{d'g}$:
\begin{align}
\mathcal{C}_{3}(q)&=2\,f^{c'd'a}\,\mathcal{D}^{c'}_{q}\,\big[U^{d'b}(z)-U^{d'b}(x)\big]\rho^{b}(x)\,\big[U^{ac}(w)-U^{ac}(y)\big]\rho^{c}(y)\notag\\
&\quad+f^{c'd'a}\,\mathcal{D}^{c'}_{q}\Big\{U^{ac}(w)U^{d'b}(z)-2\,U^{ac}(y)U^{d'b}(z)+U^{ac}(y)U^{d'b}(x)\Big\}\big[\rho^{c}(y),\rho^{b}(x)\big]. \label{eq:C3}
\end{align}
For commuting $\rho$ only the first line remains: three LO factors $[U({\rm gluon})-U({\rm source})]\rho({\rm source})$ tied by one $f$, matching the three factors $n(x'-y')$, $n(x-z)$, $n(y-w)$ of the kernel.

\paragraph{Step 2: longitudinal integrals.} With $\eta=p^{+}/k^{+}$, $s\equiv k\cdot(y'-z)$, $\vee\to\infty$ and \eqref{eq:plus}:
\begin{align}
\int_{0}^{\infty}\frac{d\eta}{1+\eta}\,e^{-i\eta s}&=e^{is}E_{1}(is)\equiv J_{0}(s),\notag\\
\int_{0}^{\infty}d\eta\left(\left[\frac{1}{\eta(1+\eta)}\right]_{+}+\frac{1}{1+\eta}\right)e^{-i\eta s}&=-\gamma_{E}-\log(is)=-\gamma_{E}-\log|s|-\tfrac{i\pi}{2}\,{\rm sgn}(s),\notag\\
\int_{0}^{\infty}d\eta\,e^{-i\eta s}&=\frac{-i}{s-i0} .
\end{align}

\paragraph{Step 3: result.}
\begin{align}
\frac{d^{3}N_{\rm III.III(1)+V+VI}}{d^{3}k}&=\frac{1}{(2\pi)^{3}}\,\frac{ig^{4}}{16\pi^{5}}\,\frac{1}{k^{+}}\int_{w,x',y',z,y,x}e^{-ik\cdot(y'-w)}\Bigg\{\notag\\
&\Big[J_{0}(s)\;n(y'-z)\!\cdot\!n(x'-y')\;\,n(y-w)\!\cdot\!n(x-z)\notag\\
&\quad+\big(\gamma_{E}+\log(is)\big)\;n(y'-z)\!\cdot\!n(x-z)\;\,n(y-w)\!\cdot\!n(x'-y')\Big]\mathcal{C}_{3}(y')\notag\\
&+\frac{i}{s-i0}\;n(y'-z)\!\cdot\!n(y-w)\;\Big[n(x-z)\!\cdot\!n(x'-y')\;\mathcal{C}_{3}(y')-n(x-z)\!\cdot\!n(x'-z)\;\mathcal{C}_{3}(z)\Big]\Bigg\}+{\rm c.c.} \label{eq:GIII3LO}
\end{align}

\subsection{Row III, second part}
Here $a=x-w$, $b=x-z$, $c=z-w$; the right colour factor is $R^{ab}$ and the left one $\mathcal{L}^{ab}_{q}\equiv f^{c'd'a}\,U^{d'b}(z)\,\mathcal{D}^{c'}_{q}$. With $\eta=p^{+}/k^{+}$ the long bracket of the first term is
\begin{align}
k^{+}\,[\dots]^{k'm}\equiv\mathfrak{B}^{k'm}(\eta)&=\delta^{k'm}\left[\frac{e+\eta\,g}{1+\eta}-\phi+\frac{d_{\perp}\,\eta\,\phi}{(1+\eta)^{2}}\right]-\omega(\eta)\,E^{k'm}-G^{k'm}-E^{mk'},\notag\\
\omega(\eta)&\equiv\left[\frac{1}{\eta(1+\eta)}\right]_{+}+\frac{1}{1+\eta}, \label{eq:B3}
\end{align}
(the two $\delta^{k'm}$ brackets of the file are $g-\phi$ and $e-\phi$; the last bracket is $G^{k'm}+E^{mk'}$), and the short brackets of the second and third terms are both $G^{mk'}$. Therefore
\begin{align}
\frac{d^{3}N_{\rm III.III(2)}}{d^{3}k}&=\frac{1}{(2\pi)^{3}}\,\frac{g^{4}}{8\pi^{5}}\,\frac{1}{k^{+}}\int_{w,x',y',z,x}e^{-ik\cdot(y'-w)}\int_{0}^{\infty}d\eta\;\frac{e^{-i\eta s}}{\eta B+A}\;n^{m}(y'-z)\,\Big\{n^{k'}(x'-y')\,\mathfrak{B}^{k'm}(\eta)\,\mathcal{L}^{ab}_{y'}\notag\\
&\hspace{4.5cm}-\eta\,G^{mk'}\big[n^{k'}(x'-y')\,\mathcal{L}^{ab}_{y'}-n^{k'}(x'-z)\,\mathcal{L}^{ab}_{z}\big]\Big\}R^{ab}\notag\\
&+\frac{1}{(2\pi)^{3}}\,\frac{g^{4}}{8\pi^{5}}\,\frac{1}{k^{+}}\int_{w,x',w',z,x}e^{-ik\cdot(w'-w)}\;\frac{A}{B\,(B-A)}\,\log\frac{B}{A}\;\;n^{m}(w'-z)\,G^{mk'}\,n^{k'}(x'-z)\;\mathcal{L}^{ab}_{z}\,R^{ab}\;+{\rm c.c.} \label{eq:GIII2}
\end{align}
In the last term the $p^{+}$ integral is done: $\int_{0}^{\infty}d\eta\,A/[B(1+\eta)(\eta B+A)]=A\log(B/A)/[B(B-A)]$, regular at $B=A$, and
$n^{m}(w'-z)G^{mk'}n^{k'}(x'-z)=\big[n(w'-z)\cdot\big(\tfrac{c}{c^{2}}-\tfrac{a}{2A}\big)\big]\big[b\cdot n(x'-z)\big]$.
If wanted, $\int_{0}^{\infty}d\eta\,e^{-i\eta s}/(\eta B+A)=B^{-1}e^{isA/B}E_{1}(isA/B)$.

\section{The $K\log$ sector across the groups}\label{sec:Ksector}

Six contributions have the kernel $u^{i}v^{j}K^{ij}\log[(x-z)^{2}/(x-w)^{2}]$, the same measure $\int_{w,w',x,y,z,x'}$ and phase $e^{-ik\cdot(w'-w)}$, and are of the form $X+{\rm c.c.}$ In units of
$\frac{1}{(2\pi)^{3}}\frac{ig^{4}}{8\pi^{5}k^{+}}\,u^{i}v^{j}K^{ij}\log\frac{B}{A}$ their colour factors are
\begin{center}
\renewcommand{\arraystretch}{1.3}
\begin{tabular}{p{3.2cm}|l}
contribution & colour factor\\ \hline
I.II, Eq.~\eqref{eq:RowIIcoord} & $\tfrac12\,\mathcal{D}_{0}^{a}\,\mathcal{S}^{a}$\\
I.III, Eq.~\eqref{eq:RowIIIfinal} & $\mathcal{D}_{0}^{a}\,f^{abd}U^{de}(x)[U^{bc}(z)-U^{bc}(y)]\rho^{e}(x)\rho^{c}(y)$\\
I.IV & $\mathcal{D}_{0}^{a}\,V^{be}\rho^{e}(y)\,R^{ab}$\\
II.I ($1a$; $1b$ is the c.c.) & $-\tfrac12\big[U^{ac'}(w')\{\rho^{c'}(x'),\rho^{b}(y)\}-U^{ad'}(x')V^{be'}\{\rho^{d'}(x'),\rho^{e'}(y)\}\big]R^{ab}$\\
III.I $+$ III.II & $-\big\{U^{ab'}(x')V^{bd'}\{\rho^{b'}(x'),\rho^{d'}(y)\}-U^{ac'}(w')V^{be'}\rho^{c'}(x')\rho^{e'}(y)$\\
 & $\quad-U^{ab'}(x')\rho^{b}(y)\rho^{b'}(x')\big\}R^{ab}$
\end{tabular}
\end{center}

\paragraph{Step 1: terms with $V$.} With the shorthand $\rho'\equiv\rho(x')$, $\rho_{y}\equiv\rho(y)$ and colour indices suppressed, the terms containing $V$ are
\begin{align}
\text{I.IV}:&\quad\big[U(x')-U(w')\big]\rho'\,V\rho_{y}\,R, \notag\\
\text{II.I}:&\quad+\tfrac12\,U(x')\,V\{\rho',\rho_{y}\}\,R, \notag\\
\text{III}:&\quad-U(x')\,V\{\rho',\rho_{y}\}\,R+U(w')\,V\rho'\rho_{y}\,R .
\end{align}
The $U(w')V$ terms cancel; the $U(x')V$ terms leave $\tfrac12\,U^{ab'}(x')V^{bd'}[\rho^{b'}(x'),\rho^{d'}(y)]R^{ab}$.

\paragraph{Step 2: terms without $V$ from II and III.}
\begin{align}
-\tfrac12\,U(w')\{\rho',\rho_{y}\}R+U(x')\rho_{y}\rho'R
=\mathcal{D}_{0}^{a}\,\rho^{b}(y)\,R^{ab}-\Big(U^{ab'}(x')-\tfrac12U^{ab'}(w')\Big)\big[\rho^{b'}(x'),\rho^{b}(y)\big]R^{ab}.
\end{align}

\paragraph{Step 3: add I.III.} Its $U(z)$ term equals $-\mathcal{D}_{0}^{a}R_{2}^{ab}\rho^{b}(y)$ where $R_{2}^{ab}=-f^{acd}U^{cb}(z)U^{de}(x)\rho^{e}(x)$ is the second term of $R^{ab}$. With Step 2 the $U(z)$ dependence cancels up to a commutator, and the first five rows of the table add to $\Omega_{0}+\Omega_{\rm comm}$,
\begin{align}
\Omega_{0}&=\mathcal{D}_{0}^{a}\Big\{f^{cbd}\,U^{ac}(w)\,\rho^{b}(y)\rho^{d}(x)-f^{abd}\,U^{bc}(y)\,U^{de}(x)\,\rho^{e}(x)\rho^{c}(y)\Big\},\\
\Omega_{\rm comm}&=\tfrac12\,U^{ab'}(x')\,V^{bd'}\big[\rho^{b'}(x'),\rho^{d'}(y)\big]R^{ab}
-\Big(U^{ab'}(x')-\tfrac12U^{ab'}(w')\Big)\big[\rho^{b'}(x'),\rho^{b}(y)\big]R^{ab}\notag\\
&\quad-\mathcal{D}_{0}^{a}\,f^{acd}\,U^{cb}(z)\,U^{de}(x)\,\big[\rho^{b}(y),\rho^{e}(x)\big].
\end{align}

\paragraph{Step 4: add I.II.} Relabelling indices, $\Omega_{0}=-\mathcal{D}_{0}^{a}f^{hbc}U^{ah}(w)\rho^{c}(y)\rho^{b}(x)+\mathcal{D}_{0}^{a}f^{abc}U^{bd}(x)U^{ce}(y)\rho^{d}(x)\rho^{e}(y)$, so
\begin{align}
\Omega_{0}+\tfrac12\,\mathcal{D}_{0}^{a}\mathcal{S}^{a}=\tfrac12\,\mathcal{D}_{0}^{a}\Big\{f^{hbc}\,U^{ah}(w)\big[\rho^{b}(x),\rho^{c}(y)\big]+f^{abc}\,U^{bd}(x)\,U^{ce}(y)\big[\rho^{d}(x),\rho^{e}(y)\big]\Big\}.
\end{align}

\paragraph{Result.}
\begin{align}
\boxed{\;\begin{aligned}
&\frac{d^{3}N_{K\log}}{d^{3}k}=\frac{1}{(2\pi)^{3}}\,\frac{ig^{4}}{8\pi^{5}}\,\frac{1}{k^{+}}\int_{w,w',x,y,z,x'}e^{-ik\cdot(w'-w)}\;u^{i}v^{j}K^{ij}\,\log\frac{(x-z)^{2}}{(x-w)^{2}}\;\Omega_{\rm tot}+{\rm c.c.},\\
&\Omega_{\rm tot}=\Omega_{\rm comm}+\tfrac12\,\mathcal{D}_{0}^{a}\Big\{f^{hbc}\,U^{ah}(w)\big[\rho^{b}(x),\rho^{c}(y)\big]+f^{abc}\,U^{bd}(x)\,U^{ce}(y)\big[\rho^{d}(x),\rho^{e}(y)\big]\Big\}
\end{aligned}\;} \label{eq:Omegatot}
\end{align}
$\Omega_{\rm tot}$ contains only commutators of $\rho$. \textbf{For commuting $\rho$ the six contributions cancel exactly.} If $[\rho^{a}(x),\rho^{b}(y)]=i\kappa f^{abc}\rho^{c}(x)\delta^{2}(x-y)$ with a convention-dependent $\kappa$, the last bracket becomes $i\kappa N_{c}\,[U^{ae}(w)+U^{ae}(x)]\rho^{e}(x)\,\delta^{2}(x-y)$ (using $f^{hbc}f^{bce}=N_{c}\delta^{he}$ and $f^{abc}U^{bd}U^{ce}=U^{ah}f^{hde}$), i.e.\ a local term with one $\rho$ less.

\section{The negative logarithms and reality}

\begin{itemize}
\item In Group I Row III the logarithms of the sign-indefinite combination $(x-w)^{2}-(x-z)^{2}$ cancel completely, with all poles $1/[(x-w)^{2}-(x-z)^{2}]^{1,2}$.
\item What remains is the constant in $\log[-B/A]=\log[B/A]+i\sigma\pi$. It does not cancel inside Row III, and no other row of the file contains a compensating term. With the prefactor $i$ it gives a real kernel:
\begin{align}
\frac{d^{3}N_{\rm I.III}}{d^{3}k}\Bigg|_{\pi}=-\,\sigma\,\frac{1}{(2\pi)^{3}}\,\frac{g^{4}}{8\pi^{4}}\,\frac{1}{k^{+}}\int e^{-ik\cdot(w'-w)}\;u^{i}\,v^{j}\,K^{ij}\;\mathcal{W}_{\rm III}\;+\;{\rm c.c.} \label{eq:piterm}
\end{align}
\item Every contribution of Groups I and III is $X+{\rm c.c.}$ and therefore real whatever the phase of $X$. For commuting $\rho$, where the colour factors are real, an $i\times$(real kernel) term gives $\sin[k\cdot(w'-w)]$ and the $\pi$ term gives $\cos[k\cdot(w'-w)]$:
\begin{align}
\frac{d^{3}N_{\rm I.III}}{d^{3}k}=\frac{1}{(2\pi)^{3}}\,\frac{g^{4}}{4\pi^{5}}\,\frac{1}{k^{+}}\int u^{i}v^{j}\,\mathcal{W}_{\rm III}
\Big\{\sin\!\big[k\cdot(w'-w)\big]\,\mathcal{R}^{ij}_{\rm real}-\sigma\pi\,\cos\!\big[k\cdot(w'-w)\big]\,K^{ij}\Big\},
\end{align}
with $\mathcal{R}^{ij}_{\rm real}=K^{ij}\log[B/A]-\tfrac12S^{ij}\log[c^{2}/A]+\mathcal{A}_{0}T^{ij}$.
\item Group I Row II has no explicit $i$, but its kernel $\mathcal{R}^{i}$ is real and odd under $(q,p)\to(-q,-p)$, so its Fourier transform is $i\times$real, consistent with \eqref{eq:RowIIcoord}.
\item Group II: $N_{1b}=(N_{1a})^{*}$; the second part is the square \eqref{eq:square}; Row IV has a real symmetric tensor and $\mathcal{H}(-s)=\mathcal{H}(s)^{*}$. All real, all logarithms positive.
\end{itemize}

\section{The simplified finite term}\label{sec:final}

After the reductions the finite term is the sum of the following pieces.
\begin{center}
\renewcommand{\arraystretch}{1.35}
\begin{tabular}{l|l|l}
 & piece & equation\\ \hline
1 & $K\log$ sector: I.II, I.IV, II.I first part, III.I, III.II and the $K\log(B/A)$ part of I.III & \eqref{eq:Omegatot} (commutators only)\\
2 & rest of I.III: $\;i\sigma\pi K^{ij}-\tfrac12S^{ij}\log\frac{c^{2}}{A}+\mathcal{A}_{0}T^{ij}$, colour $\mathcal{W}_{\rm III}$ & \eqref{eq:defR}, \eqref{eq:RowIIIfinal}\\
3 & I.V first part & \eqref{eq:RowV1}\\
4 & I.V second part ($N_{2a}+N_{2b\,\rm rest}$) & \eqref{eq:RowV2}\\
5 & II.I second part ($2a$--$2f$), perfect square & \eqref{eq:square}, \eqref{eq:KIIclosed}\\
6 & II.IV & \eqref{eq:GIIIV}\\
7 & III.III first part $+$ III.V $+$ III.VI & \eqref{eq:GIII3LO}\\
8 & III.III second part & \eqref{eq:GIII2}
\end{tabular}
\end{center}
For commuting $\rho$, piece 1 vanishes and the colour factor of piece 7 is the first line of \eqref{eq:C3}.

\section{Assumptions and points to check}\label{sec:open}

\begin{enumerate}
\item \textbf{Behaviour of I.III at $w\to z$.} $\mathcal{W}_{\rm III}$ does not depend on $w$, and for $c=z-w\to0$
\begin{align}
\mathcal{R}_{\rm III}^{ij}\;\longrightarrow\;-\frac{\delta^{ij}}{2\,(z-w)^{2}}\left\{\log\frac{(z-w)^{2}}{(x-w)^{2}}-i\sigma\pi\right\}+\mathcal{O}\!\left(\frac{1}{|z-w|}\right),
\end{align}
since $K^{ij}\to\delta^{ij}/2c^{2}$, $S^{ij}\to\delta^{ij}/c^{2}$, $K^{ij}\log[B/A]=\mathcal{O}(1/|c|)$, $\mathcal{A}_{0}T^{ij}=\mathcal{O}(1/|c|)$. $\int d^{2}w\,\log(c^{2})/c^{2}$ does not converge at $w\to z$ and the phase does not help. The source is the last term of $N_{3b}$ and the $i\sigma\pi$ term. Piece 2 of Section~\ref{sec:final} is therefore not finite by itself: either a subtraction is missing from the file, or the $\Lambda\to0$ limit does not commute with the $w$ integral there.
\item \textbf{Sign $\sigma$.} Fixed by the $i\epsilon$ of the energy denominators. The reduction assumes all ``$-1$'' inside the logarithms are continued in the same way; otherwise extra multiples of $2\pi i\,M^{ij}$ appear, and those have poles at $A=B$.
\item \textbf{Row II.} I renamed the integration momentum and assumed the $x'$ integral. The equality \eqref{eq:RowIIcoord} is exact (operator level) because $\mathcal{S}^{a}$ is exactly antisymmetric; the dropped term $\mathcal{E}^{i}$ is individually not integrable at $l\to0$ and at $p\to0$, so the statement is that it cancels in the antisymmetric combination.
\item \textbf{Plus prescriptions.} The closed-form longitudinal integrals assume \eqref{eq:plus}. A different convention changes only those closed forms, not the brackets $\mathfrak{B}$, $\mathfrak{B}_{V}$.
\item \textbf{$N_{2a}$ of I.V.} The merging with $N_{2b\,\rm rest}$ uses the limits $\int^{0}_{k^{+}}dq^{+}$ as written. If $\int_{0}^{k^{+}}$ is meant, the $\bar E^{im}$ term of $N_{2a}$ has the opposite sign, the kernels differ by that term, and the $U(x)$ cancellation is incomplete.
\item \textbf{Distributions.} $\int_{0}^{\infty}d\eta\,e^{-i\eta s}=-i/(s-i0)$ is the $\vee\to\infty$ limit in the distributional sense and includes a $\pi\delta(k\cdot(y'-z))$ piece.
\item \textbf{Commutators.} $\Omega_{\rm tot}$ and the second line of \eqref{eq:C3} are left as commutators; evaluating them needs the normalisation of $[\rho,\rho]$.
\item \textbf{Convergence} of the transverse integrals in pieces 3--8 was not examined.
\end{enumerate}

\appendix

\section{Fourier transform of $K^{ij}\log(b^{2}/a^{2})$}\label{app:FT}

Conventions: $\tilde f(k)=\int d^{2}r\,e^{ik\cdot r}f(r)$, $\lambda(k)\equiv\log(4/k^{2})-2\gamma_{E}$.
\begin{align}
\widetilde{n^{i}}(k)=\frac{2\pi i\,k^{i}}{k^{2}},\qquad
\widetilde{n^{i}\log r^{2}}(k)=\frac{2\pi i\,k^{i}}{k^{2}}\,\lambda(k),\qquad
\widetilde{\log r^{2}}(k)=-\frac{4\pi}{k^{2}}\ \text{(regularised)},\qquad
\widetilde{r^{-2}}(k)=\pi\lambda(k)+{\rm const}.
\end{align}
For a product $f(a)\,g(b)\,h(a-b)$:
\begin{align}
\int d^{2}a\,d^{2}b\;e^{iq\cdot a+il\cdot b}f(a)g(b)h(a-b)=\int\frac{d^{2}k}{(2\pi)^{2}}\,\tilde f(q-k)\,\tilde g(l+k)\,\tilde h(k).
\end{align}
Write $K^{ij}\log\frac{b^{2}}{a^{2}}=T_{1}+T_{2}+T_{3}+T_{4}$ with
$T_{1}=n^{i}(c)n^{j}(b)\log\frac{b^{2}}{a^{2}}$, $T_{2}=-n^{i}(a)n^{j}(c)\log\frac{b^{2}}{a^{2}}$, $T_{3}=-n^{i}(a)n^{j}(b)\log\frac{b^{2}}{a^{2}}$, $T_{4}=\frac{\delta^{ij}}{2c^{2}}\log\frac{b^{2}}{a^{2}}$, and $p\equiv q+l$.
\begin{itemize}
\item $T_{3}$ and $T_{4}$ factorise: $\tilde T_{3}=4\pi^{2}\frac{q^{i}l^{j}}{q^{2}l^{2}}\log\frac{q^{2}}{l^{2}}$, $\ \tilde T_{4}=\frac{2\pi^{2}\delta^{ij}}{p^{2}}\log\frac{q^{2}}{l^{2}}$.
\item The $\log b^{2}$ part of $T_{1}$ and the $\log a^{2}$ part of $T_{2}$ factorise: $-4\pi^{2}\frac{q^{i}p^{j}}{q^{2}p^{2}}\lambda(p)$ and $+4\pi^{2}\frac{p^{i}l^{j}}{p^{2}l^{2}}\lambda(p)$.
\item The remaining parts are one-loop integrals,
$\int d^{2}k\,\widetilde{\log r^{2}}(k-q)\,\varphi_{1}(k)$ with $\varphi_{1}=\frac{k^{i}(k+l)^{j}}{k^{2}(k+l)^{2}}$, and
$\int d^{2}k\,\widetilde{\log r^{2}}(k+l)\,\varphi_{2}(k)$ with $\varphi_{2}=\frac{(q-k)^{i}k^{j}}{(q-k)^{2}k^{2}}$, where
\begin{align}
\int d^{2}k\,\widetilde{\log r^{2}}(k-q)\,\varphi(k)=-4\pi\int d^{2}k\,\frac{\varphi(k)-\varphi(q)\,\theta(1-|k-q|)}{(k-q)^{2}}+4\pi^{2}(\log4-2\gamma_{E})\,\varphi(q).
\end{align}
The constants $\log4-2\gamma_{E}$ cancel against those in $\lambda(p)$.
\item After contraction with $l^{j}$ the loop integrals are done with complex coordinates ($k=k_{x}+ik_{y}$, $n(k)\leftrightarrow1/\bar k$, $n(k)\cdot l=\tfrac12[l/k+\bar l/\bar k]$): the integrand is a product of a function with simple poles in $k$ and one with simple poles in $\bar k$, and after partial fractions only
$\int d^{2}k\,\frac{1}{(k-\alpha)(\bar k-\bar\beta)}=\pi\log\frac{\Lambda^{2}}{|\alpha-\beta|^{2}}$ is needed (the cutoff $\Lambda$ cancels).
\end{itemize}
The sum is
\begin{align}
l^{j}\int d^{2}a\,d^{2}b\;e^{iq\cdot a+il\cdot b}K^{ij}\log\frac{b^{2}}{a^{2}}=2\pi^{2}\log\frac{q^{2}}{l^{2}}\left[\frac{q^{i}}{q^{2}}-\frac{p^{i}}{p^{2}}-\frac{l^{2}\,q^{i}}{p^{2}q^{2}}\right]=-2\pi^{2}\,l^{2}\,\mathcal{R}^{i}(q,p).
\end{align}
The partial-fraction sum was evaluated numerically and identified with this closed form (agreement $10^{-11}$ at 300 random momenta; scale invariance checked; the partial-fraction formula for the regulated loop integral was checked against a direct two-dimensional numerical integration).

\section{Numerical checks}\label{app:checks}
\begin{center}
\renewcommand{\arraystretch}{1.2}
\begin{tabular}{p{8.2cm}|p{8.6cm}}
statement & check\\ \hline
\eqref{eq:RowIIIfinal} vs.\ the 22 original terms of I.III & 2000 random points, complex logs, $<10^{-8}$\\
$\mathcal{A}_{0}=\theta_{w}/\Delta$; Row IV bracket $=K^{ij}$ & same run\\
$R^{i}/(p^{2}q^{2})=\mathcal{R}^{i}+2\mathcal{E}^{i}$ & random momenta, $10^{-16}$\\
\eqref{eq:FTresult} & Appendix~\ref{app:FT}\\
five colour structures $=\Omega_{0}+\Omega_{\rm comm}$; $\Omega_{0}+\tfrac12\mathcal{D}_{0}\mathcal{S}$ & $SU(2)$, non-commuting matrix-valued $\rho$, $10^{-13}$\\
$\mathcal{C}_{3}(q)$, Eq.~\eqref{eq:C3} & $SU(2)$, non-commuting $\rho$, $10^{-14}$\\
colour bracket of II.I $=R'R$ & $SU(2)$\\
\eqref{eq:KIIclosed} and \eqref{eq:square} vs.\ original $2a$--$2f$ & 3000 points, $d_{\perp}=2$ and $\neq2$, $<10^{-10}$; $\eta$ integral done numerically\\
brackets $\mathfrak{B}$ \eqref{eq:B3} and $\mathfrak{B}_{V}$ \eqref{eq:BV} vs.\ original & 500 points each, $10^{-14}$\\
longitudinal integrals ($h_{0}$, Ein, $J_{0}$, $\gamma_{E}+\log(is)$, $\mathcal{H}$) & numerical quadrature, 10 digits
\end{tabular}
\end{center}

\end{document}
